% =====================================================================
% ENSEIGNEMENTS DU STAGE
%
% CHAPITRE PROPRE À LA VERSION ENSAE. Seul main_ensae.tex l'appelle : le barème
% de l'école note sur 6 l'analyse et la réflexion sur le contexte et les outils,
% demande de préciser la part d'influence du maître de stage et des tiers, et
% attend un recul sur l'expérience. Un mémoire d'actuariat n'a rien à porter de
% tout cela, et l'Institut ne le demande pas.
%
% Le texte vient de la section « Enseignements du stage » de l'ancien rapport de
% stage raccourci (exploratory/rapport_ensae/rapport_ensae.tex, \S6). Ce qui a
% changé en le transportant ici :
%   - la hiérarchie, une \section d'article devenant un \chapter de report ;
%   - LES VINGT RENVOIS, remappés un par un sur les labels du mémoire. Les deux
%     documents ne partageaient aucun label : transporté tel quel, ce chapitre
%     aurait imprimé vingt « ?? » dans le PDF déposé ;
%   - le mot « rapport » remplacé par « mémoire » là où il désigne le document.
% =====================================================================

\chapter{Enseignements du stage}
\label{ens:ch}

\section{Le cadre de la mission et ses contraintes}
\label{ens:cadre}

Un travail de modélisation conduit dans un cabinet de conseil n'a pas le même
cadre qu'un travail conduit dans un laboratoire, et trois contraintes propres à ce
cadre ont façonné le contenu du mémoire autant que les choix statistiques.

\paragraph{Deux destinataires, deux livrables.} Le travail avait deux
destinataires, le cabinet et l'école, et ce qui est remis à l'un n'est pas ce qui
est remis à l'autre. Du côté du cabinet, le stage a donné lieu à un livrable
distinct du mémoire : l'analyse actuarielle de l'édition $2026$ d'une étude de
marché sur la cyberassurance en France, co-écrite avec le tuteur et publiée par le
cabinet \citep{RapiorKaddouri2026}. Son objet est inverse de celui du mémoire,
puisqu'elle décrit un marché observé là où le mémoire construit un modèle sur des
états hypothétiques. Elle a pourtant servi au mémoire d'une façon qui n'était pas
prévue : elle constate de l'extérieur, et sous la signature de praticiens qui
observent ce marché, l'absence de traduction de l'exposition cyber en besoin de
capital réglementaire. La prémisse du mémoire cesse ainsi d'être une affirmation
de son auteur.

\paragraph{Une donnée sous licence, et ce qu'elle impose.} La base de pertes qui
porte toute la sévérité est sous licence et ne peut pas être redistribuée. Elle ne
peut donc pas être versionnée avec le code, et un lecteur qui voudrait vérifier un
chiffre ne peut pas rejouer la chaîne depuis la donnée brute. Cette contrainte
n'est pas un détail administratif : c'est elle qui a fait naître le dispositif de
vérification de l'annexe~\ref{sec:ann-harnais}. Puisque la donnée ne peut pas
circuler, ce sont les \emph{sorties} des scripts qui sont versionnées, et le
contrôle porte sur la correspondance entre ces sorties et le document. La
traçabilité remplace la reproductibilité complète, et le document dit lequel des
deux il offre.

\paragraph{Des entités réelles, et une décision d'anonymisation.} L'application à
quatre assureurs réels aurait été plus démonstrative sous leur nom. Elle est
anonymisée, et le motif n'est pas la confidentialité des sources, qui sont des
rapports publics, mais la nature de ce que le modèle produit. Il n'utilise que le
bilan, jamais l'identité, et l'état de conformité y est \emph{supposé} : nommer une
entité reviendrait donc à lui attribuer une non-conformité que rien n'observe. La
table de correspondance existe et vit dans un script, de sorte que le travail reste
vérifiable sans être diffamatoire.

\paragraph{Le gel de la calibration, décision de gestion et non de statistique.}
La calibration a été gelée à mi-parcours, et cette décision ne relève d'aucun
critère statistique. Elle répond à un problème de conduite de projet : à partir
d'un certain volume de résultats publiés, un modèle qui bouge encore invalide en
silence ce qui a déjà été écrit, et la piste d'audit se perd. Le gel fixe une
frontière nette entre la correction d'erreur, qui reste autorisée, et la
recalibration, qui ne l'est plus. Le prix en est visible au
chapitre~\ref{ch:robustesse}, où trois écarts connus sont publiés chiffrés plutôt
que corrigés.

\paragraph{Ce que ce cadre m'a appris du métier.} La moitié du travail utile n'a
pas consisté à modéliser. Elle a consisté à rendre un résultat vérifiable par
quelqu'un qui n'a pas fait le calcul, et à écrire des énoncés qui ne dépassent pas
ce que le calcul soutient. Un modèle qu'un tiers ne peut pas contrôler n'a pas de
valeur en conseil, quelle que soit sa qualité mathématique, et c'est la différence
la plus nette que j'aie observée entre un exercice académique et une mission.

\section{Les enseignements de l'ENSAE mobilisés}
\label{ens:ensae}

Le sujet a mobilisé la formation de façon assez directe, et sur plusieurs
enseignements à la fois, ce qui est la difficulté propre d'un travail de
modélisation complet.

\begin{table}[htbp]
\centering\small
\renewcommand{\arraystretch}{1.2}
\begin{tabular}{@{}>{\raggedright\arraybackslash}p{4.2cm}>{\raggedright\arraybackslash}p{7.4cm}>{\raggedright\arraybackslash}p{2.6cm}@{}}
\toprule
\textbf{Enseignement} & \textbf{Ce qu'il fournit au travail} & \textbf{Où} \\
\midrule
Statistique des valeurs extrêmes &
dépassements de seuil, ajustement de Pareto généralisé, et surtout l'indice de
queue lu comme condition d'existence des moments &
\S\ref{soc:sec:severite} \\
Modèles à variable latente gaussienne &
dépendance entre états de conformité, transposée du risque de crédit &
chap.~\ref{chap:multietats} \\
Processus stochastiques et chaînes de Markov &
dynamique des états, trajectoire de capital, processus de branchement &
\S\ref{sec:verif-proprietes}, \S\ref{sec:res-traj} \\
Statistique inférentielle et théorie des tests &
adéquation, tests de couverture, bootstrap paramétrique, puissance &
\S\ref{soc:sec:validation}, \S\ref{soc:sec:backtest} \\
Économétrie de l'identification &
identification partielle, ensemble admissible, capital en bornes &
chap.~\ref{ch:identifiabilite}, \ref{chap:partielle} \\
Économétrie de l'évaluation &
différence de différences, test placebo, bootstrap &
chap.~\ref{ch:identifiabilite} \\
Méthodes de Monte-Carlo &
lois de perte simulées, et la discipline de publier un bruit &
\S\ref{sec:allocation}, \S\ref{sec:posture-retenue} \\
Théorie des jeux coopératifs &
valeur de Shapley pour l'attribution par canal &
\S\ref{sec:allocation} \\
Assurance non-vie et Solvabilité~II &
charge de capital, quantile réglementaire, marge de risque, allocation &
chap.~\ref{chap:reglementaire}, \S\ref{sec:detenir-transferer} \\
Programmation scientifique &
une centaine de scripts reproductibles et leurs sorties versionnées &
\S\ref{sec:ann-harnais} \\
\bottomrule
\end{tabular}
\caption{Les enseignements mobilisés et l'endroit du mémoire où chacun sert.}
\label{ens:tab:ensae}
\end{table}

\textbf{Lecture.} L'apport le plus décisif n'est pas celui que j'attendais. C'est
l'économétrie de l'identification : reconnaître qu'un paramètre n'est \emph{pas
identifié}, plutôt que constater qu'il est mal estimé, est un réflexe qui vient de
ce cours, et il a changé la nature du mémoire. La différence est considérable en
pratique, puisqu'un paramètre mal estimé appelle plus de données quand un paramètre
non identifié appelle une autre méthode. Vient ensuite la lecture de l'indice de
queue comme condition d'existence des moments, qui a permis de rejeter deux
variantes par un argument d'existence plutôt que de degré.

Deux compétences non mathématiques ont compté autant. La programmation
scientifique, parce qu'une centaine de scripts qui doivent tous rester
reproductibles pendant des mois est un problème d'ingénierie et non de calcul. Et
la rédaction, parce que la moitié du travail effectif a consisté à écrire des
énoncés qui ne dépassent pas ce que le calcul soutient.

\section{La part de l'encadrement et des tiers}
\label{ens:encadrement}

Plusieurs résultats de ce travail ne seraient pas là sans une remarque extérieure,
et cette section les recense. L'encadrement a été double : un point hebdomadaire
avec le maître de stage, préparé sous forme d'un jeu de diapositives et donnant lieu
à des demandes précises, et des échanges plus espacés avec l'encadrement académique,
dont les remarques ont porté sur le traitement de l'incertitude. Le tableau
ci-dessous distingue les deux.

\begin{table}[htbp]
\centering\small
\renewcommand{\arraystretch}{1.2}
\begin{tabular}{@{}>{\raggedright\arraybackslash}p{5.4cm}>{\raggedright\arraybackslash}p{8.4cm}@{}}
\toprule
\textbf{Demande ou remarque} & \textbf{Ce qu'elle a changé dans le travail} \\
\midrule
\multicolumn{2}{@{}l}{\emph{Du maître de stage, au point hebdomadaire}} \\
Retirer les conclusions tirées du bootstrap de l'indice de queue, un intervalle
trop large ne portant pas de conclusion &
l'intervalle reste affiché comme un fait, la lecture qu'on en faisait disparaît ;
et le retrait a fait tomber une valeur fausse d'une diapositive antérieure \\
Soupçon de double comptage dans la colonne de fermeture &
soupçon infondé, la somme des fermetures étant une identité ; une ligne
\og{}somme\fg{} sous une colonne non additive invitait cette lecture, et elle a
été retirée \\
Demande d'une table complète des leviers &
la table des trois lectures du \S\ref{sec:chiffre-de-tete} et la découverte de
l'effet de fermeture, qui change le chiffre qu'un plan de remédiation doit citer \\
Deux mots mal employés, \og{}portage\fg{} et \og{}plancher\fg{} &
le sens d'une borne était inversé, dans le texte, dans un script et dans un titre
de figure ; corrigé aux trois endroits \\
Traiter la moyenne de queue &
mesure de queue en diagnostic, et l'argument d'existence qui interdit de la
substituer au quantile \\
\midrule
\multicolumn{2}{@{}l}{\emph{De l'encadrement académique}} \\
\og{}Le quantile d'un objet incertain est un mauvais estimateur\fg{} &
le quantile prédictif et la bande de modèle ; résultat contraire à l'attente,
puisque c'est la largeur et non le point qui était fragile \\
Sensibilité aux probabilités de propagation &
le \S\ref{sec:sensibilite-propagation}, dont le résultat va contre l'inquiétude
qui la motivait \\
Niveau de la bande reportée, $90$ ou $95\,\%$ &
le niveau de $90\,\%$ est conservé, et la question a fait apparaître une
contradiction interne du mémoire (\S\ref{sec:ann-intervalle95}) \\
\bottomrule
\end{tabular}
\caption{Les demandes de l'encadrement et leur effet sur le travail.}
\label{ens:tab:encadrement}
\end{table}

\textbf{Lecture.} Aucune de ces demandes ne portait sur un calcul, et toutes ont
fait tomber un défaut. C'est le constat que je retiens de l'encadrement : la valeur
d'une relecture extérieure ne tient pas à sa compétence technique sur le détail,
elle tient à ce qu'elle interroge un énoncé hors du contexte qui le justifiait.

L'orientation du traitement des processus auto-excités mérite une mention
particulière, parce que le domaine de recherche de l'encadrement académique est
précisément celui-là. Le rejet de cet outil a été testé contre \emph{leurs}
variantes de noyau et non contre un modèle de paille, et il s'est reformulé en
cours de route en équivalence observationnelle suivie d'un choix de parcimonie
(chapitre~\ref{ch:robustesse}). Cette reformulation est plus faible en apparence et
plus honnête, et c'est celle qui tient.

\paragraph{Les tiers.} Deux sollicitations extérieures ont été lancées. Un collègue
du cabinet a été sollicité pour un second codage en aveugle du corpus de rapports
post-incident, destiné à mesurer la fiabilité entre codeurs. Ce codage n'était pas
reçu à la date de rédaction, et le dispositif d'accueil est prêt et testé sans
produire aucune donnée en son absence.

Trois praticiens du secteur ont par ailleurs été sollicités pour une relecture
critique des affirmations qualitatives du modèle, et une consultante de Nexialog a
répondu. Sa contribution n'est pas de celles qu'on espère et elle est plus utile
que celle qu'on espérait : elle n'a contredit aucune des sept affirmations, mais
son exemple d'enchaînement admet deux lectures opposées, ce qui illustre la
frontière d'identifiabilité au lieu de la lever
(chapitre~\ref{ch:identifiabilite}). Elle a aussi qualifié un arc du modèle en
distinguant l'occurrence d'un incident chez un prestataire, qu'une bonne
gouvernance ne change pas, et la maîtrise de ses conséquences, qu'elle change.
Cette réserve est déclarée et non traitée : la traiter reviendrait à déplacer cet
arc d'un canal à l'autre, donc à recalibrer, ce que le gel interdit. La répondante
n'est pas nommée. Le second codage n'est pas le préalable du résultat mais son
renforcement, et ce qu'il apporterait est une mesure d'intersubjectivité, non une
correction de niveau. Le détail de cette relecture figure au
\S\ref{sec:relecture-praticien}.

\vspace{4pt}
% SOURCES-SCRIPTS: 46 08h 85

\paragraph{L'usage d'outils d'intelligence artificielle générative.} Des outils d'intelligence
artificielle générative ont été utilisés tout au long du stage, comme assistance à la
programmation, à la vérification et à la rédaction, sous mon contrôle. Ils n'ont dispensé d'aucune
vérification : chaque nombre publié est confronté aux sorties versionnées des scripts qui le
calculent, par le dispositif décrit au \S\ref{sec:ann-harnais}. Les choix de modélisation, les
arbitrages et l'interprétation des résultats relèvent de ma responsabilité et ont été discutés
avec l'encadrement.

\section{Ce que j'en retiens}
\label{ens:retiens}

% HARNAIS-HORS-SECTION: les seuls nombres de cette section sont les comptes du
% dispositif de verification lui-meme (taux de confirmation et couverture). Les
% faire verifier par ce dispositif serait circulaire : meme declaration que la
% section D.1 du memoire, qui porte le dispositif.

\paragraph{Écrire le commentaire après avoir lu la sortie, jamais en même temps que
le code qui la produit.} C'est la leçon la plus coûteuse et je l'ai apprise sept
fois. Sept fois j'ai rédigé une conclusion pendant que le calcul tournait, et sept
fois la sortie l'a démentie : un effet annoncé convexe qui était concave, un écart
annoncé en train de se refermer qui grossissait, un ratio annoncé en effondrement
qui montait, une dépendance annoncée conservatrice que le contrôle suivant a niée,
et un compromis biais-variance que la colonne d'écart-type a exactement renversé.
Le mécanisme est toujours le même : on écrit ce que l'on attend, et le texte survit
à la mesure parce que personne ne le relit contre elle.

\paragraph{La couverture passe avant le taux.} Un indicateur de qualité qui se
calcule sur un périmètre choisi mesure le périmètre autant que la qualité. Le
dispositif de vérification affichait $99{,}6\,\%$ de nombres confirmés sur un
périmètre qui ne couvrait que $74{,}7\,\%$ des nombres publiés, et les valeurs
périmées vivaient toutes dans le quart restant. Un taux se règle en retirant une
citation, une couverture non. Cette asymétrie me paraît générale et vaut bien
au-delà de ce projet.

\paragraph{La figure est le meilleur détecteur de défauts.} Les erreurs trouvées
pendant ce stage ne l'ont presque jamais été en relisant le texte. Elles l'ont été
en regardant une figure ou une diapositive et en allant vérifier ce qu'elle
affirmait : un seuil contradictoire, six postures affichées dont deux n'étaient
définies nulle part, une même grandeur à deux valeurs, un facteur calculé à la main
et jamais imprimé, et l'anomalie de bilan du \S\ref{sec:entites-reelles}, qu'un
panneau croissant a signalée avant que je la cherche. Une figure force à énoncer un
chiffre hors du contexte qui le justifiait, et c'est là qu'un défaut devient
visible.

\section{Ce que je ferais autrement}
\label{ens:autrement}

Deux choses. J'aurais mis le dispositif de vérification en place au début et non au
milieu : construit après coup, il a fallu rattacher chaque section du document à
ses scripts, et c'est ce rattrapage qui a révélé que le chapitre central n'était
contrôlé par rien. Et j'aurais gelé la calibration plus tôt. Le gel a été décidé à
mi-parcours, après plusieurs semaines où des recalibrations partielles faisaient
bouger des nombres déjà écrits, ce qui est la façon la plus efficace de perdre une
piste d'audit.
